\documentclass[11pt]{article}

\usepackage[margin=1in]{geometry}
\usepackage[T1]{fontenc}
\usepackage[utf8]{inputenc}
\usepackage{lmodern}
\usepackage{microtype}
\usepackage{amsmath,amssymb,amsthm}
\usepackage{mathtools}
\usepackage{hyperref}
\usepackage{booktabs}
\usepackage{enumitem}
\usepackage{listings}
\usepackage{xcolor}

\hypersetup{
  colorlinks=true,
  linkcolor=blue!50!black,
  citecolor=blue!50!black,
  urlcolor=blue!50!black
}

\lstdefinelanguage{Lean4}{
  morekeywords={
    import,namespace,end,open,section,noncomputable,abbrev,
    def,theorem,lemma,where,by,fun,let,in,have,show,exact,
    intro,apply,simpa,calc
  },
  sensitive=true,
  morecomment=[l]{--},
  morecomment=[s]{/-}{-/},
  morestring=[b]"
}

\lstset{
  language=Lean4,
  basicstyle=\ttfamily\small,
  keywordstyle=\bfseries\color{blue!50!black},
  commentstyle=\itshape\color{green!40!black},
  stringstyle=\color{red!50!black},
  showstringspaces=false,
  breaklines=true,
  frame=single,
  framerule=0.3pt,
  xleftmargin=0.5em,
  xrightmargin=0.5em,
  keepspaces=true,
  columns=fullflexible,
  literate=
    {¬}{{$\neg$}}1
    {∃}{{$\exists$}}1
    {→}{{$\to$}}1
    {≤}{{$\leq$}}1
}

\newtheorem{theorem}{Theorem}[section]
\newtheorem{lemma}[theorem]{Lemma}
\theoremstyle{remark}
\newtheorem{remark}[theorem]{Remark}

\title{Navier--Stokes Grassroots:\\
Ground-Up Background Theory for the Goertzel SG-Cole-Hopf Program}
\author{}
\date{\today}

\begin{document}
\maketitle

\begin{abstract}
This note records the \emph{grassroots} side of the current work around Ben
Goertzel's December 11, 2025 Navier--Stokes manuscript.  It is intentionally
not the blocker note.  The companion crux-first analysis lives separately in
\texttt{ns\_crux.tex}.

The goal here is narrower and more constructive: isolate reusable background
theory that would still matter even if the SG-Cole-Hopf route needs revision.
The first machine-checked kernel now formalizes the exact algebraic inequalities
behind the manuscript's identity-energy route: finite carré-du-champ energy,
log-gradient reduction under a positive lower bound, the \(-2\nu \log\) scaling
factor, and the Cauchy--Schwarz push-down from coefficient energy at the
identity to pointwise vorticity/frame bounds.  It now also includes an abstract
Cole--Hopf identity interface: if one can supply uniform positivity, an
identity-energy bound for \(P_t\phi\), and a frame-curl bound, the whole
algebraic chain to a uniform vorticity estimate is derived in Lean.  And it now
adds a local kernel-founded semigroup shell on top of Mettapedia's own
Markov-kernel infrastructure, so the next NS layer can stay inside the local
library rather than reaching outward prematurely.  The newest pass also adds a
positive approximation shell: if a future finite-mode or truncation theorem
shows convergence of the manuscript's radius and statistic observables, then
the induced cutoff potential and Cole--Hopf field convergence now follow
automatically, and any eventual absolute truncation bound on those quantities
passes to the limit as well.  On the identity-energy side, a new finite-mode
limit shell now says that eventual truncation positivity and eventual
truncation coefficient-energy bounds promote to honest limiting
\texttt{ColeHopfIdentityData}, so the existing vorticity bridge can be reused on
the limit object.  And the same finite-mode transfer layer now fuses with the
local kernel-semigroup shell, so the kernel-founded NS interface can also be
recovered directly on the limit data.  The newest pass packages the whole
manuscript-shaped truncation target in one Lean structure: chart-observable
convergence, cutoff/Cole--Hopf convergence, finite-mode positivity and
identity-energy transfer, and kernel-founded vorticity consequences now sit in
one explicit object.  And the newest concrete step lands the D.4 chart side in
an actual finite-mode model: a squared weighted chart radius, a bounded linear
observable, the chart-cutoff datum \(W_0 = \chi \cdot \ell\), the heat datum
\(\phi = \exp(-W_0/(2\nu))\), explicit \(W_0/\phi\) bounds, and truncation
convergence for both \(W_0\) and \(\phi\) are now machine-checked in one local
module.  The latest pass then packages the first genuinely Cole--Hopf-shaped
shared truncation instance into direct named kernel and vorticity objects, so
the cutoff-convergence and uniform vorticity consequences can now be cited from
one concrete package instead of from the whole abstraction stack.  The newest
step then replaces the static selected-coefficient \(d\Phi\) surrogate by a
nonnegative-time heat-decayed one, while preserving the same shared
truncation, positivity, and coefficient-energy consequences.

This does \emph{not} prove 3D Navier--Stokes regularity.  What it does is make
the positive architecture cleaner: if the hard analytic bridges can be supplied,
the identity-energy argument now lands in a small Lean-checked interface rather
than in manuscript prose alone.
\end{abstract}

\tableofcontents

\section{Scope}

The December 2025 manuscript
\emph{``Global Regularity of 3D Navier--Stokes via SG-Cole-Hopf Program''}
proposes a route with four main layers:
\begin{enumerate}[leftmargin=1.5em]
  \item a polynomially weighted divergence-free Fourier frame on the torus;
  \item an SG diffusion and Cole--Hopf transform on the diffeomorphism group;
  \item an identity-based push-down back to Eulerian velocity fields;
  \item a first-order energy estimate at the identity strong enough to control
  vorticity and close the BKM criterion.
\end{enumerate}

The companion crux note asks where this route may fail.  This note asks a
different question:
\begin{quote}
What mathematical infrastructure do we want in Lean even if the final
regularity proof still needs serious repair?
\end{quote}

That changes the work style.  We prefer small exact interfaces over speculative
closure.  In particular:
\begin{enumerate}[leftmargin=1.5em]
  \item separate algebraic identities from heavy infinite-dimensional analysis;
  \item formalize the inequalities the manuscript repeatedly reuses;
  \item make the hard analytic inputs explicit, so future repairs can plug into
  them rather than hiding them.
\end{enumerate}

There is also a build-sanity rule now:
\begin{quote}
build from Mettapedia plus the locally pinned \texttt{v4.28.0} mathlib first,
and treat outside Lean repos as reference material only unless a direct local
need appears.
\end{quote}

\section{Existing Base}

The grassroots route is not starting from zero.  The project already has:
\begin{center}
\begin{tabular}{p{0.46\textwidth}p{0.42\textwidth}}
\toprule
\textbf{Module} & \textbf{Grassroots role} \\
\midrule
\texttt{Mathlib} finite sums / real inequalities & Cauchy--Schwarz and square-root algebra \\
\shortstack[l]{\texttt{Mettapedia/ProbabilityTheory/}\\\texttt{MarkovCategory/Kernels.lean}} & stable local kernel semantics \\
\shortstack[l]{\texttt{Mettapedia/FluidDynamics/}\\\texttt{NavierStokes/IdentityEnergy.lean}} & new NS identity-energy kernel \\
\shortstack[l]{\texttt{Mettapedia/FluidDynamics/}\\\texttt{NavierStokes/ColeHopfInterface.lean}} & abstract manuscript-shaped identity bridge \\
\shortstack[l]{\texttt{Mettapedia/FluidDynamics/}\\\texttt{NavierStokes/KernelSemigroupInterface.lean}} & local kernel-founded semigroup shell \\
\shortstack[l]{\texttt{Mettapedia/FluidDynamics/}\\\texttt{NavierStokes/CutoffContinuity.lean}} & generic cutoff continuity / instability shell \\
\shortstack[l]{\texttt{Mettapedia/FluidDynamics/}\\\texttt{NavierStokes/ColeHopfTopology.lean}} & local \(W \mapsto \phi\) continuity / instability shell \\
\shortstack[l]{\texttt{Mettapedia/FluidDynamics/}\\\texttt{NavierStokes/ApproximationInterface.lean}} & positive truncation / approximation shell \\
\shortstack[l]{\texttt{Mettapedia/FluidDynamics/}\\\texttt{NavierStokes/ManuscriptChartDatum.lean}} & concrete D.4-style chart datum on a finite-mode model \\
\shortstack[l]{\texttt{Mettapedia/FluidDynamics/}\\\texttt{NavierStokes/ManuscriptChartPackage.lean}} & concrete chart datum plugged into the combined truncation package \\
\shortstack[l]{\texttt{Mettapedia/FluidDynamics/}\\\texttt{NavierStokes/SharedApproximationPackage.lean}} & one shared `approx` for chart and identity lanes \\
\shortstack[l]{\texttt{Mettapedia/FluidDynamics/}\\\texttt{NavierStokes/GeometricSharedApproximation.lean}} & first explicit shared truncation instance with positivity and energy bounds \\
\shortstack[l]{\texttt{Mettapedia/FluidDynamics/}\\\texttt{NavierStokes/GeometricColeHopfSharedApproximation.lean}} & Cole--Hopf-shaped shared instance driven by the actual chart heat datum \\
\shortstack[l]{\texttt{Mettapedia/FluidDynamics/}\\\texttt{NavierStokes/GeometricColeHopfPackage.lean}} & direct concrete package-level kernel and vorticity consequences for that instance \\
\shortstack[l]{\texttt{Mettapedia/FluidDynamics/}\\\texttt{NavierStokes/GeometricColeHopfHeatApproximation.lean}} & nonnegative-time heat-decayed \(d\Phi\) model on the same shared truncation family \\
\shortstack[l]{\texttt{Mettapedia/FluidDynamics/}\\\texttt{NavierStokes/FiniteModeLimit.lean}} & finite-mode positivity / energy transfer to the limit \\
\shortstack[l]{\texttt{Mettapedia/FluidDynamics/}\\\texttt{NavierStokes/FiniteModeKernelLimit.lean}} & finite-mode limit transfer into the local kernel shell \\
\shortstack[l]{\texttt{Mettapedia/FluidDynamics/}\\\texttt{NavierStokes/ManuscriptTruncationPackage.lean}} & combined manuscript-shaped truncation target \\
\texttt{ns\_crux.tex} & companion blocker note, kept separate on purpose \\
\bottomrule
\end{tabular}
\end{center}

So the honest answer to ``how much can we build on Mettapedia?'' is:
\begin{quote}
quite a lot for the finite algebra and interface design, not much yet for the
heavy PDE/stochastic geometry itself.
\end{quote}

The current NS-specific Lean files are deliberately modest.  They do not attempt
Kunita theory, infinite-dimensional geometry, or PDE well-posedness.  They
extract the clean finite algebra, the exact hypothesis interface that the
manuscript's identity-only argument depends on, and now a local kernel shell
for future semigroup-facing work.

\section{The New Lean Kernel}

The new module
\texttt{Mettapedia/FluidDynamics/NavierStokes/IdentityEnergy.lean}
defines a finite carré-du-champ style energy
\[
  \Gamma(D) := \sum_i D_i^2
\]
for a finite family of directional derivatives or coefficients.

\subsection{Log-Energy Reduction}

The first kernel theorem is the exact positivity estimate used in the manuscript
when passing from \(\Phi\) to \(\log \Phi\).

\begin{theorem}[Finite log-gradient bound]
If \(m > 0\) and \(m \le F\), then
\[
  \Gamma(D/F) \le \Gamma(D) / m^2.
\]
\end{theorem}

\noindent
This is formalized as \texttt{gamma\_div\_le}.  It is elementary, but it
extracts one of the manuscript's recurrent steps into a reusable statement with
no stochastic or geometric baggage.

The second theorem packages the specific Cole--Hopf scaling used in the NS
route.

\begin{theorem}[The \(-2\nu \log\) energy factor]
If \(m > 0\) and \(m \le \Phi\), then
\[
  \Gamma\!\left((-2\nu)\,D/\Phi\right)
  \le \frac{4\nu^2}{m^2}\,\Gamma(D).
\]
\end{theorem}

\noindent
This is \texttt{gamma\_negTwoNuLog\_le}.  It is exactly the algebraic part of
the passage from a bound on \(\Gamma(P_t \phi)(\mathrm{id})\) to a bound on
\(\Gamma(W(t))(\mathrm{id})\) once \(W(t) = -2\nu \log(P_t \phi)\) and a
positive lower bound for \(P_t \phi\) are available.

\subsection{Frame Push-Down and Vorticity}

The second half of the file formalizes the finite Cauchy--Schwarz step behind
the vorticity estimate.  Given coefficients \(a_i\) and a frame \(E_i(x)\), define
\[
  \mathrm{frameEval}(a,E,x) := \sum_i a_i E_i(x).
\]

\begin{theorem}[Finite frame Cauchy--Schwarz]
For every \(x\),
\[
  |\mathrm{frameEval}(a,E,x)|
  \le \sqrt{\Gamma(a)}\;\sqrt{\Gamma(E(\cdot,x))}.
\]
\end{theorem}

\noindent
This appears in Lean as \texttt{abs\_frameEval\_le}.  The specialized theorem
\texttt{abs\_vorticity\_le} is the manuscript-shaped version where the frame
components are \(\mathrm{curl}\,E_i(x)\).

So the current Lean kernel already isolates the exact positive shape:
\begin{quote}
identity coefficient energy + uniform frame-curl energy
\(\Longrightarrow\) pointwise vorticity bound.
\end{quote}

That is not the whole Navier--Stokes proof.  It is the algebraic spine that the
hard analytic work has to feed.

\subsection{The Cole--Hopf Identity Interface}

The second new module,
\texttt{Mettapedia/FluidDynamics/NavierStokes/ColeHopfInterface.lean},
packages the manuscript-shaped hypotheses into one structure:
\begin{itemize}[leftmargin=1.5em]
  \item a positive profile \(\Phi(t)\),
  \item directional derivatives \(D_i\Phi(t)\) at the identity,
  \item a uniform lower bound \(m_\phi\),
  \item a uniform identity-energy bound on \(\Gamma(D\Phi(t))\),
  \item a curl frame with a uniform frame-energy bound.
\end{itemize}

From that interface, Lean now derives the full algebraic bridge:
\begin{enumerate}[leftmargin=1.5em]
  \item \texttt{gamma\_Wcoeff\_le}: a uniform bound on the coefficients of
  \(W(t) = -2\nu\log\Phi(t)\);
  \item \texttt{abs\_vorticity\_le}: a pointwise vorticity bound at any time and
  point;
  \item \texttt{abs\_vorticity\_le\_uniform}: the uniform version over all times
  and points.
\end{enumerate}

This is exactly the kind of reuse Mettapedia is already good at: small abstract
interfaces with derived consequences, leaving the deep analytic hypotheses
explicit rather than implicit.

\subsection{A Local Kernel-Founded Semigroup Shell}

The next new module,
\texttt{Mettapedia/FluidDynamics/NavierStokes/KernelSemigroupInterface.lean},
does not try to prove SG diffusion facts.  It does something more conservative:
it re-expresses the NS evolution layer in terms of Mettapedia's own stable
kernel semantics from
\texttt{Mettapedia/ProbabilityTheory/MarkovCategory/Kernels.lean}.

Concretely, the file introduces a minimal local structure:
\begin{itemize}[leftmargin=1.5em]
  \item a measurable state object \(\Omega\),
  \item time-indexed endokernels \(P_t : \Omega \rightsquigarrow \Omega\),
  \item identity and composition laws for the kernel family,
  \item the same identity-only Cole--Hopf fields already used in the algebraic
  interface.
\end{itemize}

Lean then proves that all existing coefficient and vorticity consequences are
recovered unchanged from this kernel-founded package.  So the current NS lane
is now grounded on a local semigroup shell without pretending that the hard
analytic semigroup estimates are already done.

\subsection{A Generic Cutoff Continuity Repair Shell}

The newest local helper,
\texttt{Mettapedia/FluidDynamics/NavierStokes/CutoffContinuity.lean},
extracts the exact mechanism behind the chart-cutoff issue without committing
to any SG-specific geometry.

It defines a generic cutoff potential
\[
  W(x) = \eta(r(x))\,s(x)
\]
and proves two small but useful facts:
\begin{enumerate}[leftmargin=1.5em]
  \item if the cutoff, radius observable, and statistic are continuous, then
  \(W\) is continuous;
  \item if a convergent sequence has the cutoff factor tending to the wrong
  limit, then \(W\) cannot be continuous at the limit point.
\end{enumerate}

This helps the route constructively.  The grassroots side can now say:
\begin{quote}
any repair of the chart-cutoff step must supply continuity of the actual radius
observable in the chosen topology, or replace the cutoff construction entirely.
\end{quote}

That does not solve the crux.  It does turn the crux into a local reusable
interface theorem instead of leaving it as manuscript-level warning prose.

\subsection[A Local W-to-phi Topology Shell]{A Local \(W \mapsto \phi\) Topology Shell}

The next new module,
\texttt{Mettapedia/FluidDynamics/NavierStokes/ColeHopfTopology.lean},
pushes the same idea one step further.  It formalizes the real Cole--Hopf map
\[
  \phi = \exp\!\left(-\frac{W}{2\nu}\right)
\]
as a local topological wrapper over the generic cutoff potential.

Its role is again deliberately modest:
\begin{enumerate}[leftmargin=1.5em]
  \item prove that if \(W\) is continuous, then the induced \(\phi\) is
  continuous;
  \item prove that if a sequence forces the wrong limit for \(W\), then it also
  forces the wrong limit for \(\phi\), provided \(\nu \neq 0\);
  \item specialize this to the chart-cutoff potential
  \(W(x)=\eta(r(x))\,s(x)\).
\end{enumerate}

This matters because the crux note's instability is not only about \(W_0\).  It
propagates through the manuscript's actual Cole--Hopf transform.  The new file
turns that propagation into a local reusable theorem rather than leaving it as
an informal remark.

\subsection{A Positive Approximation Repair Shell}

The newest module,
\texttt{Mettapedia/FluidDynamics/NavierStokes/ApproximationInterface.lean},
records the constructive side of the same issue.

It introduces a tiny structure
\texttt{CutoffApproximationData(radius, statistic)}
whose content is exactly this:
\begin{itemize}[leftmargin=1.5em]
  \item an approximant family \(x \mapsto x_n\),
  \item convergence of the radius observable \(r(x_n) \to r(x)\),
  \item convergence of the statistic \(s(x_n) \to s(x)\).
\end{itemize}

From those hypotheses Lean now derives:
\begin{enumerate}[leftmargin=1.5em]
  \item convergence of the cutoff factor \(\eta(r(x_n)) \to \eta(r(x))\);
  \item convergence of the cutoff potential
  \(W(x_n)=\eta(r(x_n))\,s(x_n)\to W(x)\);
  \item convergence of the corresponding Cole--Hopf field
  \(\phi(x_n) \to \phi(x)\);
  \item inheritance of any eventual absolute truncation bound for \(W(x_n)\) or
  \(\phi(x_n)\) by the limit object.
\end{enumerate}

This is the clean positive counterpart to the crux note's mismatch argument.
It says the route does not need vague reassurance about truncations.  It needs
a theorem of one specific shape:
\begin{quote}
the manuscript's actual finite-mode approximants must force convergence of the
radius and statistic observables in the topology being used.
\end{quote}

If that theorem exists, the cutoff/Cole--Hopf passage is no longer handwavy.
The new Lean shell is where such a repair theorem plugs in.  It also means that
``finite-mode first, infinite-mode later'' is no longer just a strategy slogan:
once the finite truncations have a uniform bound and the relevant observables
converge, the same bound transfers to the limit inside Lean.

\subsection{A Concrete D.4 Chart Datum}

The new module
\texttt{Mettapedia/FluidDynamics/NavierStokes/ManuscriptChartDatum.lean}
turns the manuscript's D.4 formula shape into an actual finite-mode Lean model
instead of leaving it only as an interface.

It starts from the existing geometric mode state and defines:
\begin{enumerate}[leftmargin=1.5em]
  \item a squared weighted chart radius \(r^2\), closer to the manuscript's
  \(\|\kappa(g)\|_{H^m}^2\);
  \item a bounded linear observable \(\ell\) with explicit envelope control;
  \item the concrete chart datum \(W_0 = \chi(r^2)\,\ell\);
  \item the concrete heat datum
  \(\phi = \exp(-W_0/(2\nu))\).
\end{enumerate}

Lean then proves three kinds of facts.
\begin{enumerate}[leftmargin=1.5em]
  \item Finite-mode truncation convergence: the squared radius, the bounded
  observable, \(W_0\), and \(\phi\) all converge along truncations.
  \item Boundedness of the observable and the chart datum: once the cutoff
  \(\chi\) is bounded, the resulting \(W_0\) is bounded explicitly.
  \item Lower and upper bounds for \(\phi\): the local D.4 positivity envelope
  is now derived inside Lean from the \(W_0\) bound.
\end{enumerate}

This is still not the manuscript's full SG chart.  But it is a meaningful
upgrade in fidelity.  The chart side is no longer only ``some abstract radius
and statistic.''  There is now a concrete local model where the D.4
construction is correct, bounded, and truncation-stable.

The follow-up file
\texttt{Mettapedia/FluidDynamics/NavierStokes/ManuscriptChartPackage.lean}
then plugs this concrete D.4 chart model straight into the already-existing
\texttt{ManuscriptTruncationPackage.lean}.  So the chart side is not only
concrete in isolation; it now lands directly in the combined package that also
carries the kernel-founded vorticity consequences.

The next new file,
\texttt{Mettapedia/FluidDynamics/NavierStokes/SharedApproximationPackage.lean},
then removes the last interface duplication.  It records one shared
\(\mathrm{approx} : \mathbb{N} \to G \to G\) and derives from that same family:
\begin{enumerate}[leftmargin=1.5em]
  \item the chart-side cutoff/Cole--Hopf approximation package;
  \item the finite-mode identity-energy package;
  \item the combined manuscript-shaped truncation package.
\end{enumerate}

So the current grassroots architecture is now cleaner:
\begin{quote}
one approximation family, one chart side, one identity side, one combined
package.
\end{quote}

The next file,
\texttt{Mettapedia/FluidDynamics/NavierStokes/GeometricSharedApproximation.lean},
is the first actual instance of that architecture.  It takes the concrete D.4
chart datum from the geometric mode state, uses the same truncation family
\texttt{truncateModes}, reads \(d\Phi\) off finitely many selected coefficients
of the same truncated state, and proves:
\begin{enumerate}[leftmargin=1.5em]
  \item eventual positivity for the same \(\Phi\)-family;
  \item eventual finite \(\Gamma(d\Phi)\) bounds for the same family;
  \item direct inhabitation of the full
  \texttt{SharedApproximationPackage.lean}.
\end{enumerate}

So the current grassroots lane no longer says only
``if someone supplies a shared truncation family, we know what to do.''  It now
says:
\begin{quote}
here is one concrete shared family, and here are the positivity and energy
bounds for that same family.
\end{quote}

The next file,
\texttt{Mettapedia/FluidDynamics/NavierStokes/GeometricColeHopfSharedApproximation.lean},
pushes one step closer to the manuscript.  It replaces the toy positive-offset
\(\Phi\) with the actual chart heat datum
\[
  \phi = \exp\!\left(-\frac{W_0}{2\nu}\right)
\]
from the D.4-style chart model, and it defines \(d\Phi\) as that same heat
datum multiplied by finitely selected coefficients of the same truncated state.

That matters because the positivity and energy bounds are now proved for a
family whose \(\Phi\)-side is genuinely Cole--Hopf-shaped.  Concretely, Lean
now proves for the same shared truncation family:
\begin{enumerate}[leftmargin=1.5em]
  \item eventual lower positivity of \(\phi\) from the explicit \(W_0\) bound;
  \item eventual finite \(\Gamma(d\Phi)\) bounds for the \(\phi\)-scaled
  selected-coefficient model;
  \item direct inhabitation of the full
  \texttt{SharedApproximationPackage.lean} by this more faithful choice.
\end{enumerate}

So the situation is now better than before:
\begin{quote}
we no longer have only a shared truncation toy model; we also have a shared
Cole--Hopf-shaped heat-datum model.
\end{quote}

The follow-up file
\texttt{Mettapedia/FluidDynamics/NavierStokes/GeometricColeHopfPackage.lean}
then turns that instance into named concrete package objects.  It exposes:
\begin{enumerate}[leftmargin=1.5em]
  \item the concrete manuscript-shaped truncation package for the geometric
  Cole--Hopf instance;
  \item the associated concrete finite-mode kernel data and kernel semigroup
  data;
  \item the named vorticity field for that package;
  \item direct theorems for cutoff-potential convergence, Cole--Hopf
  convergence, coefficient-energy control, and uniform vorticity bounds.
\end{enumerate}

That is a real usability step.  Later formalization no longer needs to unpack
the whole shared-approximation tower just to cite the first concrete
Cole--Hopf-shaped NS package.

The newest file,
\texttt{Mettapedia/FluidDynamics/NavierStokes/GeometricColeHopfHeatApproximation.lean},
then takes one more real step toward the manuscript's evolution story without
pretending to have proved the full SG semigroup.  It keeps the same D.4 chart
heat datum \(\phi\), keeps the same shared truncation family, but replaces the
static selected-coefficient identity model by
\[
  d\Phi_i(t) = \phi \, e^{-t\,k_i}\, a_{k_i}
\]
for nonnegative time and selected modes \(k_i\).  Lean proves:
\begin{enumerate}[leftmargin=1.5em]
  \item truncation convergence for the heat-decayed coefficient model;
  \item a uniform finite \(\Gamma(d\Phi(t))\) bound, since the heat factors are
  bounded by \(1\) for \(t \ge 0\);
  \item direct inhabitation of the same
  \texttt{SharedApproximationPackage.lean} by this more genuinely
  time-dependent family.
\end{enumerate}

So the identity side is no longer only
\begin{quote}
static selected coefficients multiplied by \(\phi\).
\end{quote}
It is now
\begin{quote}
selected coefficients with an explicit nonnegative-time heat decay, still tied
to the same concrete chart datum and the same finite-mode truncation family.
\end{quote}

\subsection{A Finite-Mode Identity-Energy Limit Shell}

The next new module,
\texttt{Mettapedia/FluidDynamics/NavierStokes/FiniteModeLimit.lean},
does the same kind of cleanup for the manuscript's identity-energy lane.

It proves three things:
\begin{enumerate}[leftmargin=1.5em]
  \item pointwise convergence of a finite coefficient family implies convergence
  of its finite carré-du-champ energy \(\Gamma\);
  \item eventual finite-mode energy bounds \(\Gamma(D_n)\le A\) pass to the
  limit coefficients;
  \item eventual positivity lower bounds \(m \le \Phi_n\) also pass to the
  limit profile.
\end{enumerate}

Those ingredients are then packaged into a structure
\texttt{FiniteModeColeHopfData} whose theorem
\texttt{toColeHopfIdentityData} produces the exact existing
\texttt{ColeHopfIdentityData} interface on the limit object.  So once a future
analytic argument supplies:
\begin{quote}
finite-mode convergence, eventual lower positivity, and eventual coefficient
energy bounds,
\end{quote}
the already verified vorticity bridge applies to the limit automatically.

This is the manuscript's ``finite-mode first, infinite-mode later'' pattern in
a sharper form than before.  It no longer stops at convergence of observables.
It also transfers the actual identity-energy hypotheses needed by the current
Lean vorticity route.

The follow-up module
\texttt{Mettapedia/FluidDynamics/NavierStokes/FiniteModeKernelLimit.lean}
then combines this limit-transfer shell with the already local
\texttt{KernelSemigroupInterface.lean}.  So if a future analytic argument
supplies both a kernel semigroup and the finite-mode convergence/bound package,
the present kernel-founded vorticity estimate is available on the limit object
with no further algebra.

Finally,
\texttt{Mettapedia/FluidDynamics/NavierStokes/ManuscriptTruncationPackage.lean}
packages the full current target in one place.  It combines:
\begin{enumerate}[leftmargin=1.5em]
  \item chart-observable convergence for the cutoff side;
  \item cutoff-potential and Cole--Hopf convergence consequences;
  \item finite-mode positivity and identity-energy transfer to the limit;
  \item the resulting kernel-founded vorticity bound on the limit data.
\end{enumerate}

That matters because the next analytic theorem no longer has to be described
vaguely as ``something feeding several shells.''  It now has one exact Lean
landing point.

\section{Why This Matters}

The value of this kernel is not that it proves a Millennium problem.  The value
is that it separates the route into two layers.

\begin{enumerate}[leftmargin=1.5em]
  \item \textbf{Reusable algebraic layer.}
  Once one has a lower bound on \(P_t \phi\), an identity-energy bound for
  \(\Gamma(P_t\phi)(\mathrm{id})\), and a frame-curl constant, the rest of the
  vorticity estimate is just these finite inequalities.
  \item \textbf{Hard analytic input layer.}
  The real burden is now explicit: produce those bounds in the actual infinite-
  dimensional SG setting, justify the semigroup/flow passage, and justify the
  push-down theorem analytically.
\end{enumerate}

That split is precisely what a grassroots program should want.  If the SG route
survives, these lemmas will still be needed.  If the SG route mutates, the same
identity-energy interface may survive even then.

\section{Positive Manuscript Infrastructure}

Even on a charitable read, not every part of the manuscript deserves the same
confidence.  But the positive architecture is real enough to build around.

\subsection{Polynomial Fourier Frame}

The manuscript gives a concrete divergence-free Fourier frame with polynomial
orbit weights on \(T^3\).  This is not just inspirational prose.  It is a
specific normalization scheme designed to make the relevant \(\ell^2\)-type
sums finite for \(m \ge 5\).  Grassroots formalization can treat this as a real
mathematical object, separate from the later stochastic and geometric layers.

\subsection{Finite-Mode First, Infinite-Mode Later}

The manuscript repeatedly follows a sensible proof pattern:
\begin{quote}
prove the estimate for finite truncations first, then pass to the full
infinite-mode object.
\end{quote}

That is exactly the right shape for formalization.  It suggests that future Lean
work should not start with the full diffeomorphism-group diffusion.  It should
start with finite coefficient families, finite frame sums, and explicit
truncation interfaces.

\subsection{Schwartz Exhaustion Is Background, Not Center Stage}

On a closer read, the manuscript's periodization and torus-exhaustion appendix
is more substantial than a first skim suggests.  For grassroots purposes, the
important point is structural:
\begin{quote}
the \(\mathbb{R}^3\) step mostly wants uniform control of the same torus-side
constants already needed earlier.
\end{quote}

So the best constructive work remains the torus-side background theory, not an
early jump to the Euclidean limit.

\section{What Still Needs Heavy Analysis}

The new Lean kernel also clarifies what it does \emph{not} solve.

\begin{enumerate}[leftmargin=1.5em]
  \item It does not prove the local propagation estimate for
  \(\Gamma(P_t \phi)(\mathrm{id})\).
  \item It does not prove the manuscript-shaped truncation theorem that would
  feed convergence of the radius and statistic observables into the new
  approximation shell, or the eventual positivity / coefficient-energy bounds
  into the new finite-mode limit shell.
  \item It does not justify the manuscript's push-down theorem at a fully
  analytic infinite-dimensional level.
  \item It does not close BKM by itself; it only packages the algebraic part of
  that closure once the analytic hypotheses are supplied.
\end{enumerate}

That honesty is a feature, not a weakness.  The point of grassroots work is to
shrink the unknown region and expose the remaining burden sharply.

\section{Next Grassroots Targets}

The next constructive steps are fairly clear.

\begin{enumerate}[leftmargin=1.5em]
  \item Instantiate the new shared approximation package with a genuinely
  manuscript-shaped truncation family: the chart-side D.4 datum and the
  finite-mode identity-energy datum should come from the same explicit
  truncation family, with convergence of the actual radius/statistic
  observables and eventual lower positivity / coefficient-energy bounds for the
  identity data.  That now feeds directly into
  \texttt{ApproximationInterface.lean} and
  \texttt{FiniteModeLimit.lean}, or equivalently into the combined
  \texttt{ManuscriptTruncationPackage.lean} through
  \texttt{SharedApproximationPackage.lean}.  The new geometric and geometric
  Cole--Hopf instances are now explicit realizations of this pattern, and the
  new heat-decayed Cole--Hopf instance is a more time-aware one, but the next
  goal is still a more faithful manuscript-shaped package in which the chosen
  \(d\Phi\) coefficients are justified by the actual SG/identity differential
  story rather than by a finite selected-coefficient heat surrogate.
  \item Strengthen the new local kernel-founded shell from ``kernel semigroup +
  identity interface'' to ``kernel semigroup + positivity/energy propagation''.
  That is the next honest place where a real analytic theorem can plug in.
  \item Formalize finite-frame Sobolev-to-vorticity constants cleanly, so the
  current \texttt{abs\_vorticity\_le} theorem can be plugged into an explicit
  frame family rather than an abstract coefficient system.
\end{enumerate}

The reason for this order is straightforward: land a small exact kernel before
spending effort on a speculative infinite-dimensional closure.

\section{Conclusion}

The SG-Cole-Hopf route is not formally closed.  But it is already structured
enough that its algebraic core can be extracted and verified independently.

That is what the current grassroots pass accomplishes.  The new Lean module does
not prove global regularity, but it does prove that the manuscript's
identity-energy route has a small reusable spine:
\[
  \Gamma(P_t \phi)(\mathrm{id})
  \Longrightarrow
  \Gamma(W(t))(\mathrm{id})
  \Longrightarrow
  |\omega(t,x)|
\]
once the hard analytic hypotheses are fed in.

That is exactly the kind of background theory we want from the ground up.

\end{document}
